\section{Methodology}

We present \textbf{Flux \& Key}, a training-free, reference-guided image editing pipeline for
multi-modal diffusion transformers. Given a source image
$\mathbf{I}_A \in \mathbb{R}^{H \times W \times 3}$, a reference image
$\mathbf{I}_B \in \mathbb{R}^{H' \times W' \times 3}$ containing the object of interest, a
binary spatial mask $\mathbf{M}_A \in \{0,1\}^{H \times W}$ delineating the edit region in
$\mathbf{I}_A$, and a text description $\mathcal{T}$, the goal is to produce an edited image
$\hat{\mathbf{I}}$ such that:
(i)~the visual identity of the specific object in $\mathbf{I}_B$ is faithfully transferred into
the masked region;
(ii)~all pixels outside $\mathbf{M}_A$ are reconstructed with high fidelity; and
(iii)~no model fine-tuning or iterative ODE inversion is required.
The pipeline consists of five stages detailed below.

% -----------------------------------------------------------------------
\subsection{Reference Object Extraction}
% -----------------------------------------------------------------------

We first isolate the target object in $\mathbf{I}_B$ using a grounded segmentation module.
Grounding DINO~\cite{liu2023grounding} localises the object region via text-conditioned
bounding-box prediction; SAM~2~\cite{ravi2024sam2} then produces a precise binary object
mask $\mathbf{M}_B$.
The reference crop is obtained as
\begin{equation}
    \mathbf{I}_B^{\text{crop}} = \mathbf{I}_B \odot \mathbf{M}_B,
    \label{eq:crop}
\end{equation}
where $\odot$ denotes element-wise multiplication.
To address the scale mismatch between the reference object and the target region we resize
$\mathbf{I}_B^{\text{crop}}$ to match the bounding-box extent of $\mathbf{M}_A$:
\begin{equation}
    \mathbf{I}_B^{\text{aligned}} =
        \operatorname{Resize}\!\left(\mathbf{I}_B^{\text{crop}},\;
                                      \operatorname{bbox}(\mathbf{M}_A)\right).
    \label{eq:resize}
\end{equation}
Critically, rather than encoding $\mathbf{I}_B^{\text{aligned}}$ through a CLIP
encoder---which is optimised for semantic similarity and discards high-frequency instance
detail~\cite{radford2021clip}---we pass it directly through FLUX's VAE encoder $\mathcal{E}$,
preserving the full spatial and spectral structure of the reference:
\begin{equation}
    \mathbf{z}_B = \mathcal{E}\!\left(\mathbf{I}_B^{\text{aligned}}\right)
        \in \mathbb{R}^{h_A \times w_A \times C},
    \label{eq:ref_encode}
\end{equation}
where $h_A, w_A$ are the spatial dimensions of the edit-region latent.

% -----------------------------------------------------------------------
\subsection{Inversion-Free Background Preservation via Frequency Locking}
% -----------------------------------------------------------------------

Rather than caching background key-value pairs from a costly DDIM inversion
pass~\cite{zhu2025kvedit}, we preserve the background through a single clean forward pass,
maintaining the pipeline's inversion-free character throughout.

\paragraph{Source feature extraction.}
We compute the clean latent $\mathbf{z}_A = \mathcal{E}(\mathbf{I}_A)$ and perform one forward
pass through the frozen FLUX transformer $f_\theta$ at $t=0$ to extract source attention
features at each block $\ell$:
\begin{equation}
    \left\{\mathbf{K}_{\text{src}}^{(\ell)},\;
           \mathbf{V}_{\text{src}}^{(\ell)}\right\}_{\ell=1}^{L}
    = f_\theta(\mathbf{z}_A).
    \label{eq:src_features}
\end{equation}

\paragraph{Frequency-locked background attention.}
During denoising, at each timestep and block $\ell$, let
$\mathbf{K}_{\text{tgt}}^{(\ell)}, \mathbf{V}_{\text{tgt}}^{(\ell)}$
be the keys and values produced from the current noisy latent.
We define complementary high- and low-pass frequency filters
\begin{equation}
    \mathcal{H}(\hat{\mathbf{X}})_{u,v} =
        \hat{\mathbf{X}}_{u,v}\cdot\mathbf{1}\!\left[\|(u,v)\| > \gamma\right],
    \qquad
    \mathcal{L}(\hat{\mathbf{X}})_{u,v} =
        \hat{\mathbf{X}}_{u,v}\cdot\mathbf{1}\!\left[\|(u,v)\| \leq \gamma\right],
    \label{eq:filters}
\end{equation}
where $\mathcal{F}$ denotes the 2-D Discrete Fourier Transform and $\gamma$ is the
frequency-split threshold.
For background tokens $\mathcal{B} = \{i : \mathbf{M}_A^{(i)} = 0\}$ we construct
\begin{align}
    \mathbf{K}_{\mathcal{B}}^{(\ell)}
        &= \mathcal{F}^{-1}\!\left(
              \mathcal{H}\!\left(\mathcal{F}\!\left(\mathbf{K}_{\text{src}}^{(\ell)}\right)\right)
            + \mathcal{L}\!\left(\mathcal{F}\!\left(\mathbf{K}_{\text{tgt}}^{(\ell)}\right)\right)
           \right)\!\bigg|_{\mathcal{B}},
    \label{eq:bg_K} \\
    \mathbf{V}_{\mathcal{B}}^{(\ell)}
        &= \mathcal{F}^{-1}\!\left(
              \mathcal{H}\!\left(\mathcal{F}\!\left(\mathbf{V}_{\text{src}}^{(\ell)}\right)\right)
            + \mathcal{L}\!\left(\mathcal{F}\!\left(\mathbf{V}_{\text{tgt}}^{(\ell)}\right)\right)
           \right)\!\bigg|_{\mathcal{B}}.
    \label{eq:bg_V}
\end{align}

High-frequency components of the \emph{source} keys and values encode the structural layout,
edge geometry, and fine texture of the background, locking these signals in place without any
inversion.
Low-frequency components from the \emph{target} branch carry global illumination and semantic
context, allowing the background to subtly adapt its lighting when a visually different object
is inserted---a property that hard KV-freezing~\cite{zhu2025kvedit} suppresses but that
photorealistic compositing requires.
Unlike FIA-Edit~\cite{yang2026fia}, which applies the same frequency fusion over the entire
image without a mask, our formulation confines the lock exclusively to $\mathcal{B}$ tokens,
eliminating edit-content bleed into the background by construction.

% -----------------------------------------------------------------------
\subsection{Selective Stochastic Inversion}
% -----------------------------------------------------------------------

Full DDIM inversion applies the reverse ODE trajectory to the entire image, accumulating
numerical error that corrupts even unedited regions~\cite{wallace2023edict}.
Inversion-free methods avoid this cost but give the model insufficient freedom to regenerate
a structurally different object in the foreground.
We resolve this tension with \textbf{Selective Stochastic Inversion (SSI)}, which applies
controlled forward-direction noise exclusively to masked tokens.

Let $\mathbf{z}_A^{\mathcal{F}}$ denote the foreground token subset of $\mathbf{z}_A$,
where $\mathcal{F} = \{i : \mathbf{M}_A^{(i)} = 1\}$.
SSI initialises the denoising trajectory for the edit region as
\begin{equation}
    \tilde{\mathbf{z}}_0^{\mathcal{F}}
        = \sqrt{1-\tau^2}\;\mathbf{z}_A^{\mathcal{F}} + \tau\,\boldsymbol{\varepsilon},
    \qquad \boldsymbol{\varepsilon} \sim \mathcal{N}(\mathbf{0}, \mathbf{I}),
    \label{eq:ssi}
\end{equation}
where $\tau \in [0, 1]$ is the \emph{stochasticity coefficient} governing edit aggressiveness.
Background tokens are initialised from the clean source latent,
$\tilde{\mathbf{z}}_0^{\mathcal{B}} = \mathbf{z}_A^{\mathcal{B}}$, and are never perturbed.
Unlike FSI-Edit~\cite{yang2026fsi}, which injects noise globally and risks content bleeding
beyond the edit boundary, SSI is spatially confined to $\mathcal{F}$, eliminating
mask-leakage artefacts by construction.

The stochasticity coefficient $\tau$ is calibrated per edit category: for tasks requiring
wholesale object replacement (e.g., change object, add object),
$\tau \in [0.7, 0.9]$ affords sufficient regeneration freedom; for attribute-level
modifications (e.g., change colour, change texture),
$\tau \in [0.3, 0.5]$ preserves more source structure.

% -----------------------------------------------------------------------
\subsection{Spectral-Aware Dual-Source KV Injection}
% -----------------------------------------------------------------------

The central architectural contribution is a mechanism for injecting reference appearance into
the foreground attention stream while preserving the spatial layout prescribed by the text
prompt.
At transformer block $\ell$, let
$\mathbf{K}_{\text{text}}^{(\ell)}, \mathbf{V}_{\text{text}}^{(\ell)}$ denote the keys and
values produced from the text branch, and
$\mathbf{K}_{\text{ref}}^{(\ell)}, \mathbf{V}_{\text{ref}}^{(\ell)}$ those obtained by
projecting the reference latent $\mathbf{z}_B$ through the same QKV weight matrices.
We fuse the two sources in the frequency domain:
\begin{align}
    \tilde{\mathbf{K}}^{(\ell)}
        &= \mathcal{F}^{-1}\!\left(
              \alpha^{(\ell)}\,
              \mathcal{F}\!\left(\mathbf{K}_{\text{text}}^{(\ell)}\right)
            + \left(1 - \alpha^{(\ell)}\right)
              \mathcal{F}\!\left(\mathbf{K}_{\text{ref}}^{(\ell)}\right)
           \right),
    \label{eq:fg_K} \\
    \tilde{\mathbf{V}}^{(\ell)}
        &= \mathcal{F}^{-1}\!\left(
              \beta^{(\ell)}\,
              \mathcal{F}\!\left(\mathbf{V}_{\text{text}}^{(\ell)}\right)
            + \left(1 - \beta^{(\ell)}\right)
              \mathcal{F}\!\left(\mathbf{V}_{\text{ref}}^{(\ell)}\right)
           \right),
    \label{eq:fg_V}
\end{align}
where $\alpha^{(\ell)}, \beta^{(\ell)} \in [0,1]$ are layer-wise mixing coefficients.

\paragraph{Frequency decomposition rationale.}
The low-frequency spectrum of $\mathbf{K}_{\text{text}}$ encodes coarse spatial layout and
semantic structure---information that governs scale, pose, and placement of the inserted
object.
The high-frequency spectrum of $\mathbf{K}_{\text{ref}}$ encodes edge sharpness, surface
texture, and fine geometric detail that distinguish one physical instance from another.
Naive convex combination in feature space conflates these orthogonal signals; FFT mixing
decouples them, allowing the text branch to guide layout while the reference branch
contributes appearance fidelity.

\paragraph{Layer-wise scheduling.}
We adopt a linearly decaying schedule for $\alpha^{(\ell)}$:
\begin{equation}
    \alpha^{(\ell)}
        = \alpha_{\max} - \frac{\ell}{L}
          \left(\alpha_{\max} - \alpha_{\min}\right),
    \label{eq:schedule}
\end{equation}
with $\alpha_{\max} = 0.8$ and $\alpha_{\min} = 0.2$.
Early blocks (small $\ell$) are text-dominant, preserving global structural layout; late
blocks (large $\ell$) are reference-dominant, sharpening instance-specific appearance.
The schedule for $\beta^{(\ell)}$ mirrors Eq.~\eqref{eq:schedule}.

The full per-block attention for foreground tokens is then
\begin{equation}
    \mathbf{A}_{\mathcal{F}}^{(\ell)}
        = \operatorname{Softmax}\!\!\left(
              \frac{\mathbf{Q}_{\mathcal{F}}^{(\ell)}\,
                    \tilde{\mathbf{K}}^{(\ell)\top}}{\sqrt{d}}
          \right)\tilde{\mathbf{V}}^{(\ell)}.
    \label{eq:fg_attn}
\end{equation}
Background tokens attend exclusively to
$\{\mathbf{K}_{\mathcal{B}}^{(\ell)}, \mathbf{V}_{\mathcal{B}}^{(\ell)}\}$
from Eqs.~\eqref{eq:bg_K}--\eqref{eq:bg_V}, and the two streams are never mixed.

% -----------------------------------------------------------------------
\subsection{Position Embedding Re-indexing}
% -----------------------------------------------------------------------

The reference crop $\mathbf{I}_B^{\text{aligned}}$ is encoded at an arbitrary absolute
position in its own coordinate frame.
To ensure spatial coherence with the target region we re-index its 2-D rotary position
embeddings (RoPE) to match the token-grid coordinates of $\mathcal{F}$ within Image~A:
\begin{equation}
    \mathbf{pe}_{\text{ref}}^{(i)}
        = \operatorname{RoPE}\!\left(
              \operatorname{row}\!\left(\mathcal{F}^{(i)}\right),\;
              \operatorname{col}\!\left(\mathcal{F}^{(i)}\right)
          \right),
    \label{eq:rope}
\end{equation}
where $\bigl(\operatorname{row}(\mathcal{F}^{(i)}), \operatorname{col}(\mathcal{F}^{(i)})\bigr)$
are the 2-D grid indices of the $i$-th foreground token in Image~A's full latent.
This resolves spatial inconsistencies arising from scale mismatch without any explicit affine
transformation.

% -----------------------------------------------------------------------
\subsection{Inference and Post-Processing}
% -----------------------------------------------------------------------

The complete denoising loop proceeds as follows.
At each timestep~$t$ and transformer block~$\ell$:
\begin{enumerate}
    \item Background tokens attend to frequency-locked
          $\{\mathbf{K}_{\mathcal{B}}^{(\ell)}, \mathbf{V}_{\mathcal{B}}^{(\ell)}\}$
          (Eqs.~\eqref{eq:bg_K}--\eqref{eq:bg_V}).
    \item Foreground tokens attend to spectral-mixed
          $\{\tilde{\mathbf{K}}^{(\ell)}, \tilde{\mathbf{V}}^{(\ell)}\}$
          (Eqs.~\eqref{eq:fg_K}--\eqref{eq:fg_V}).
    \item The denoised latent is decoded by the VAE decoder $\mathcal{D}$ at the final step.
\end{enumerate}

\paragraph{Lightweight harmonisation.}
A final post-processing step aligns the colour statistics of the decoded edit region with the
local neighbourhood of $\mathbf{M}_A$ in $\mathbf{I}_A$ via histogram matching, correcting
residual lighting inconsistencies introduced by the domain gap between
$\mathbf{I}_A$ and $\mathbf{I}_B$:
\begin{equation}
    \hat{\mathbf{I}}^{\mathcal{F}}
        = \operatorname{HistMatch}\!\left(
              \mathcal{D}(\mathbf{z}_{\text{final}})^{\mathcal{F}},\;
              \mathbf{I}_A^{\mathcal{N}(\mathbf{M}_A)}
          \right),
    \label{eq:histmatch}
\end{equation}
where $\mathbf{I}_A^{\mathcal{N}(\mathbf{M}_A)}$ denotes the pixels in $\mathbf{I}_A$
within a neighbourhood of the mask boundary.

\paragraph{Computational complexity.}
The method requires two forward passes (one on $\mathbf{I}_A$, one on
$\mathbf{I}_B^{\text{aligned}}$) plus $T$ denoising steps each with two FFT operations per
block, each of complexity $\mathcal{O}(N \log N)$ where $N = |\mathcal{F}|$.
No ODE inversion is performed at any stage.
The overall pipeline is summarised in Figure~\ref{fig:pipeline}.
